%Tema para beamer "CCM4" versión 1
%Desarrollo por Erick David Luna Núñez y 
%Fernanda Barajas Hernandez

\documentclass{beamer}

\usepackage[utf8]{inputenc}
\usepackage{heuristica}
\usepackage[T1]{fontenc}
\usepackage[heuristica,vvarbb,bigdelims]{newtxmath}
\renewcommand*\oldstylenums[1]{\textosf{#1}}
\usepackage[sfdefault,scaled=.85]{FiraSans}
\usepackage{graphicx}
\usepackage[spanish]{babel} 
\usepackage[pages=some]{background}
\pagenumbering{arabic}



%%Se define el "environment" teorema
\newtheorem{teorema}{Teorema}

%Definir el autor con el estilo definido (el textbf y el uso del color son herramientas del diseño, no necesario borrar)
\title{ \textbf{Nombre de} \\ {\color{mostazaccm} la plática}}

%Nombre del autor
\author{Nombre del autor} 

%Fecha o evento en que se presentará la plática
\date{Congreso muy importante 2023} 

%correo del expositor o incluir posibles colaboradores
\institute{fulano@matmor.unam.mx} 


%%Tema de beamer "CCM-4"
\usetheme{ccm4}

\begin{document}
%Define el fondo de la primer diapositiva
{\setbeamertemplate{background}{
\includegraphics[width=\the\paperwidth,height=\the\paperheight]{images/P4.png}}
%lo anterior es para definir el fondo de la primer diapositiva

\begin{frame}
  \titlepage %Necesario para generar la portada
\end{frame}

} %aquí termina el cambio de fondo

\begin{frame}
\tableofcontents %Imprime la tabla de contenido
\end{frame}

\section{Introducción} %%Título de la sección (Opcional)
\subsection{Ejemplo}
\begin{frame}
  \frametitle{Título de la diapositiva}
  \framesubtitle{Subtítulo} %%Subtítulo de la diapositiva (opcional)
  Algo de texto en la parte superior
  
  \[
   B * \pi =\int_{0}^{\infty} \frac{5x^2}{\sqrt{a+b}}\, dx
   \]

  \begin{itemize}
    \item[\checkmark] Un elemento en la lista %%[\checkmark] muestra una palomita al inicio de la línea
    \item Segundo elemento
    \item Otro elemento
    \item Y finalmente, otro más
  \end{itemize}
\end{frame}

\section{Sección de ejemplo 1}
\begin{frame}{Otra diapositiva} %%Otra forma (más corta) de poner el título a la diapositiva
Texto adicional. No tiene un propósito en particular más que ocupar algo de espacio.
Texto adicional. No tiene un propósito en particular más que ocupar algo de espacio.
Texto adicional. No tiene un propósito en particular más que ocupar algo de espacio.
Texto adicional. No tiene un propósito en particular más que ocupar algo de espacio.
\( \sum_0^{\infty} a_i \) 
Texto adicional. No tiene un propósito en particular más que ocupar algo de espacio.
\end{frame}

\section{Sección de ejemplo 2} %%Otra sección
\begin{frame}
    \frametitle{Un teorema}
    
    \begin{teorema} %%Uso del "environment" definido al inicio del documento.
        Los números primos son infinitos
    \end{teorema}

    Y algo de texto
    
    \centering \includegraphics[width=0.45 \paperwidth]{images/Manifold1.png}
\end{frame}

\section{Teorema con referencia}
\begin{frame}{Ejemplo}
\framesubtitle{Incluye referencia}
    \begin{teorema}[Teorema de Stokes] 
        Sea $M$ una variedad compacta orientada con frontera, $\omega$ una $k-1$-forma sobre $M$, entonces $$\int_{M} d\omega = \int_{\partial M} \omega$$
    \end{teorema}

    la prueba se puede revisar en: \cite{spivak2018calculus}
\end{frame}

\section{Referencias}
\begin{frame}{Referencias}
    %posibles estilos para bibliografía: unsrt , siam , plain , ieeetr , alpha , acm , abbrv
    \bibliographystyle{apalike}
  \bibliography{references}
\end{frame}


\begin{frame}{Gracias}
    
\end{frame}

\end{document}
